% =====================================================================
%  A WORKED REPRODUCTION of
%  Miki Aoyagi (2023), "Consideration on the learning efficiency of
%  multiple-layered neural networks with linear units" (preprint,
%  ssrn 4404877; submitted to Neural Networks).
%
%  Purpose of THIS document (read first):
%   - This is NOT a copy. It is a re-derivation: every statement is
%     restated in fixed notation, every proof is reproduced with the
%     gaps filled in by us, and commentary is interleaved.
%   - It is the controller's reference for the `aoyagi-full` expedition:
%     the place where we (a) hold the vision of Aoyagi's actual method,
%     (b) work out the paper's typos / underspecifications, and
%     (c) record the bridge to the Lean formalisation.
%   - Commentary is tagged so it can be skimmed or stripped:
%       \fnote{...}  formalisation note (Lean target / Mathlib hook)
%       \tflag{...}  TYPO or UNDERSPECIFICATION in the paper + our fix
%       \gnote{...}  geometric reformulation (codim S(t)=Mval bridge)
%       \qmark{...}  open question / to-verify (exact-arithmetic check)
%
%  Status: all sections filled.
%   [x] 0  Notation + reading guide
%   [x] 1  Model, loss, RLCT framework (Def 1, Lemma 1, Def 2)
%   [x] 2  The two theorems (Thm 1 = RRR/L=2, Thm 2 = general-L) + Def 3
%   [x] 3.1 Theorem 3  (product reduction / iterated block elimination)
%   [x] 3.2 Theorem 4  (reduction to the deepest singular point)
%   [x] 3.3 The recursive blow-up (Cases 1 & 2) — THE MOUNTAIN
%   [x] 3.4 The candidate thresholds M_{s,k} and the min
%   [x] 4  Lemmas 3–5 (arithmetic minimisation + order theta)
%   [x] 5  The RRR (L=2) specialisation, packaged
%   [x] 6  Typos / underspecifications, collected
%
%  2026-07-21 ELDER ANNOTATION PASS (strike-wave generation folded in):
%   - "Lean status" \fnote{}s added at each theorem/lemma: what is now PROVEN
%     sorry-free at the expedition tip vs paper-claimed vs deferred.
%   - FULL PAGE-IMAGE VERIFICATION of pp.14--22 done this pass (poppler on root;
%     the seats lacked PDF tooling): the inductive statement + b_i closed form +
%     Jacobian ledger (p.15), Case-1 exponents (pp.16--17), Case-2 label + exponent
%     + the T-profile totality display (pp.19--20), the M_{s,k} formula + t~=0
%     read-off (p.22) are all IMAGE-CONFIRMED against this document. New ledger
%     items (T-E), (T-F); see also verify-owed-math-audit.md (the owed-math list).
% =====================================================================
\documentclass[11pt]{article}
\usepackage{amsmath,amssymb,amsthm,mathtools,bm}
\usepackage[margin=1in]{geometry}
\usepackage{xcolor}

\theoremstyle{plain}
\newtheorem{theorem}{Theorem}
\newtheorem{lemma}{Lemma}
\newtheorem{proposition}{Proposition}
\theoremstyle{definition}
\newtheorem{definition}{Definition}
\newtheorem{example}{Example}
\theoremstyle{remark}
\newtheorem*{remark}{Remark}

% commentary tags
\newcommand{\fnote}[1]{\par\noindent{\color{blue!60!black}$\vdash$ \textsf{[Formalisation]} #1}\par}
\newcommand{\tflag}[1]{\par\noindent{\color{red!70!black}$\boxtimes$ \textsf{[Typo/Underspec]} #1}\par}
\newcommand{\gnote}[1]{\par\noindent{\color{green!45!black}$\triangle$ \textsf{[Geometry]} #1}\par}
\newcommand{\qmark}[1]{\par\noindent{\color{orange!80!black}$?$ \textsf{[To verify]} #1}\par}

% notation
\newcommand{\R}{\mathbb{R}}
\newcommand{\rk}{\operatorname{rank}}
\newcommand{\codim}{\operatorname{codim}}
\newcommand{\norm}[1]{\lVert #1 \rVert}
\newcommand{\ideal}[1]{\langle #1 \rangle}
\newcommand{\rlct}{\lambda}          % log canonical threshold / learning coefficient
\newcommand{\rorder}{\theta}         % its order (multiplicity of the pole)
\newcommand{\Mval}{\mathrm{Mval}}
\newcommand{\Ms}{M^{\ast}}           % the integer M in Def 3 (renamed to avoid clash with widths)

\title{A worked reproduction of Aoyagi (2023):\\
the learning coefficient of deep linear networks}
\author{(\texttt{aoyagi-full} expedition --- controller's worked copy)}
\date{}

\begin{document}
\maketitle

\tableofcontents

% =====================================================================
\section{Notation and reading guide}
% =====================================================================

We fix notation once, deviating from the paper only to remove clashes.

\begin{itemize}
\item \textbf{Depth.} There are $L$ weight matrices $A^{(1)},\dots,A^{(L)}$ (an
``$L$-layer linear network''; the paper's \emph{three-layered} network is $L=2$).
Layer widths are $H^{(1)},\dots,H^{(L+1)}$ with $A^{(s)}$ of size
$H^{(s)}\times H^{(s+1)}$, so the product
$\prod_{s=1}^{L}A^{(s)} := A^{(1)}A^{(2)}\cdots A^{(L)}$ is $H^{(1)}\times H^{(L+1)}$.
\item \textbf{Data.} input $X\in\R^{H^{(L+1)}}$, output $Y\in\R^{H^{(1)}}$,
$Y=\prod_s A^{(s)}X+\text{noise}$.
\item \textbf{True parameter} $\bar w=\{\bar A^{(s)}\}$; $r:=\rk\!\big(\prod_s\bar A^{(s)}\big)$.
\item \textbf{Reduced widths} $M^{(s)}:=H^{(s)}-r$, for $s=1,\dots,L+1$.
\item $E_s$ is the $s\times s$ identity; $O$ a zero block of the size demanded by context.
\item For a matrix $C$, $\norm{C}=\sqrt{\sum_{ij}c_{ij}^2}$ (Frobenius) and
$\ideal{C}$ is the ideal generated by all entries $c_{ij}$.
\end{itemize}

\gnote{Throughout, the single object that controls everything is the
\emph{geometry of the zero fibre} $\{\prod_s C^{(s)}=0\}$ of the reduced core.
The paper computes $\rlct$ by an explicit resolution; we will keep a parallel
``geometric'' reading in which $\rlct=\tfrac12\min_t \codim S(t)$, where $S(t)$
ranges over the nested-rank strata. The two must agree; their agreement is the
spine of the formalisation, and the geometric form is what lets us bypass the
paper's Definition~3 (see \S2).}

% =====================================================================
\section{Model, loss, and the RLCT framework}
% =====================================================================

\subsection{The statistical model and the optimal set}

The model density is
\[
 p(Y\mid X,w)=\tfrac{1}{(\sqrt{2\pi})^{H^{(1)}}}
   \exp\!\Big(-\tfrac12\,\norm{Y-\textstyle\prod_s A^{(s)}X}^2\Big),
 \qquad p(X,Y\mid w)=p(Y\mid X,w)\,r(X),
\]
with the true density $r(X,Y)$ obtained by plugging $\bar w$. The average log
loss $L(w)=-\mathbb E_{X,Y}[\log p]$ is minimised exactly on the \emph{optimal set}
(the fibre of the multiplication map)
\[
 W_0=\Big\{\,w:\ \textstyle\prod_s A^{(s)}=\prod_s \bar A^{(s)}\,\Big\}.
\]
\fnote{This is the expedition's \texttt{optimalSet d B} with
$B=\prod_s\bar A^{(s)}$ and \texttt{dlnLoss\,d\,B\,$w$}$=\norm{\prod_s A^{(s)}-B}^2_F$.
The Kullback function $K(w)$ of the model is (up to the Gaussian normalisation,
which is constant in $w$) equal to $\tfrac12\norm{\prod_s A^{(s)}-\prod_s\bar A^{(s)}}^2$;
its RLCT and the loss's RLCT coincide because they generate the same ideal
(Lemma~\ref{lem:ideal} below). So we may and do work directly with the
polynomial $\norm{\prod_s A^{(s)}-B}^2$.}

\subsection{The log canonical threshold (Definition 1)}

\begin{definition}[Log canonical threshold; Aoyagi Def.\ 1]\label{def:rlct}
For an analytic $F$ and a $C^\infty$ bump $\varphi$ with compact support in a
neighbourhood $U$ of $w^\ast$,
\[
 \rlct_{w^\ast}(F,\varphi)=\sup\Big\{\,c:\ \int_U |F|^{-c}\,\varphi(w)\,dw<\infty\,\Big\}
 \qquad(\text{real field}),
\]
equivalently the largest pole of the zeta function $\int_U|F|^{-z}\varphi\,dw$,
$z\in\mathbb C$; let $\rorder_{w^\ast}(F,\varphi)$ be the order of that pole.
If $\varphi(w^\ast)\neq0$ the value is independent of $\varphi$ and we write
$\rlct_{w^\ast}(F)$, $\rorder_{w^\ast}(F)$.
For an ideal $J=\ideal{F_1,\dots,F_m}$ put $\rlct_{w^\ast}(J)=\rlct_{w^\ast}(F_1^2+\cdots+F_m^2)$.
\end{definition}

\fnote{This is \texttt{rlctAt} (Rung S0). Two things must be built, not cited:
(i) the $\sup$-of-finite-$c$ definition over an open neighbourhood (a germ; the
expedition verified \texttt{rlctAt} is the germ threshold, so a fixed compact box
needs a finite-subcover step --- already built as \texttt{box\_integrable\_of\_lt\_rlctAtOn\_deepest});
(ii) $\varphi$-independence (used freely below). The \emph{order} $\rorder$ is the
analytically harder second deliverable.}

\begin{lemma}[RLCT depends only on the ideal; Aoyagi Lemma 1, after {[30,31,32]}]\label{lem:ideal}
Let $J=\ideal{F_1,\dots,F_n}$ on a neighbourhood of $w^\ast$ and let $G_1,\dots,G_m$
be analytic. If $G_1,\dots,G_m\in J$ then
$\rlct_{w^\ast}(G_1^2+\cdots+G_m^2)\ge\rlct_{w^\ast}(F_1^2+\cdots+F_n^2)$.
In particular if $J=\ideal{G_1,\dots,G_m}$ as well, the two RLCTs are equal.
\end{lemma}

\fnote{This is Rung S1, the workhorse. It is what lets every block-elimination
and change-of-variable below preserve the RLCT (they preserve the ideal). It is
also what equates the loss-RLCT with the Kullback-RLCT. \emph{Must be proven}
(monotonicity of $\rlct$ under ideal inclusion). Mathlib has no RLCT; the
inequality is the elementary direction (an inclusion of ideals gives a pointwise
domination $\sum G^2\le c\sum F^2$ locally, hence the integral comparison).}

\fnote{\textbf{Lean status (strike wave 1, 2026-07-21): PROVEN, sorry-free.}
Object A landed in \texttt{Core.Aoyagi.IdealInvariance}: the one-sided
$\le$ (\texttt{rlctAt\_sumSqFam\_le\_of\_germRepresents}), the two-sided equality
(\texttt{rlctAt\_sumSqFam\_eq\_of\_germ\_eq}), and the \emph{weighted} forms
(\texttt{wrlctAt\_*}) that the resolution change-of-variables consumes. Two
statement-level guards were REQUIRED beyond the paper's prose, both elder-adjudicated
with counterexamples: (i) the \emph{junk-$0$ guard} \texttt{LocallyNullZeros} --- Lean's
\texttt{rpow} has $0^{-c}=0$, so a positive-measure zero set of $\sum G_i^2$ would
\emph{inflate} the Lean threshold (false without the guard; dischargeable for
polynomial/analytic families, \texttt{Waypoint.locallyNullZeros\_sumSqFam\_of\_polynomial});
(ii) measurability of the weight and the dominated family (a Vitali-modulated
continuous-at-$x$-only cofactor otherwise annihilates one admissible set). Germ ideal
membership is encoded as an explicit representation (\texttt{GermRepresents}, coefficients
continuous at $x$; dom-wide \texttt{RegionRepresents} for the resolution charts).}

\gnote{\textbf{Why Lemma 1 is load-bearing, exactly (thread-27, exact algebra).} The
recursion's unit transforms $Q,P$ are \emph{non-orthogonal}, so the Frobenius norm is NOT
preserved: after the resolution the pulled-back loss is \emph{not literally} $\sum b_i^2$
--- it is a sum of squares of elements \emph{generating the same ideal}. Verified exactly
at the coupled instance $(3,3,4)$ (\texttt{g-coupled-334-diagb.py}). Without Lemma~1 the
whole diagonalisation chain computes the threshold of the wrong function; with it, only
the ideal matters. This is the precise sense in which the ideal route replaces the dead
chart route.}

\begin{definition}[Aoyagi Def.\ 2]
$\norm{C}=\sqrt{\sum_{ij}|c_{ij}|^2}$ and $\ideal{C}=\ideal{\{c_{ij}\}}$.
\end{definition}

\subsection{The resolution route to $\rlct$ (Hironaka + the monomial rule)}

Applying Hironaka's theorem to $K(w)$ yields a proper analytic map
$g:Q\to V\subseteq W$ with, in local coordinates $(u_1,\dots,u_d)$ on $Q$,
\[
 K(g(u))=u_1^{2k_1}\cdots u_d^{2k_d},\qquad
 |g'(u)|\,\varphi(g(u))=u_1^{h_1}\cdots u_d^{h_d},
\]
and then
\[
 \boxed{\ \rlct=\min_{\text{charts }U}\ \min_{1\le j\le d}\frac{h_j+1}{2k_j}\,,\qquad
 \rorder=\max_{u}\#\Big\{j:\tfrac{h_j+1}{2k_j}=\rlct\Big\}.\ }
\]
\fnote{\textbf{This boxed rule is S2 --- the \emph{only} permitted citation.}
Everything else (the construction of $g$ as explicit charts, the pullback being
monomial$\times$unit, the Jacobian being a monomial, the charts covering, and the
$\min$ over charts being realised by the joint integral) is on our side. In Lean,
S2 is \texttt{monomial\_rlct}: $\texttt{monomialThreshold}\,d\,k\,h
=\min_j (h_j+1)/(2k_j)$.}

\fnote{\textbf{Lean status (strike wave 1, 2026-07-21): the per-chart boxed rule is
now a THEOREM, not a citation.} Object C landed sorry-free
(\texttt{Core.Aoyagi.MonomialRLCT}): \texttt{monomialSumSq\_wrlctAt\_eq} proves
$\mathrm{wrlct}(W,\sum b_k^2)=\min_{d:\,k_d>0}(h_d{+}1)/(2k_d)$ for a monomial family
under the \emph{divisibility chain} $b_{k_0}\mid b_k$ (Aoyagi's own invariant, the
principal-normal-crossing hypothesis) against $W=\text{(monomial)}\times\text{(unit)}$,
and \texttt{monomialSumSq\_two\_mul\_wrlctAt\_eq\_min} gives the DLN unit-multiplicity
integer form $2\rlct=\min(h_d{+}1)$. The chain hypothesis is load-bearing: the coupled
family $(u_0u_1^2,\,u_0^2u_1)$ has axis value $1$ but true threshold $2/3$, and is
\emph{excluded} (\texttt{not\_divChain\_coupled\_example}). Guards: \texttt{Measurable
unit} (delta-ratified) $+$ the junk-$0$ guard. What remains OPEN of this boxed display
is exactly its \emph{atlas} half --- the ``$\min$ over charts $U$'' --- the sorried
frontier leaf \texttt{rlctAt\_sumSqFam\_eq\_iInf\_charts}
(\texttt{Core.Aoyagi.ProductResolution}); the per-chart value
\texttt{Chart.two\_mul\_wrlctAt\_eq\_chartMin} and the wired
\texttt{Resolution.two\_mul\_rlctAt\_eq\_divisorMin} are proven modulo that leaf.}

\gnote{Read the boxed rule with $k_j\equiv1$ (every exceptional divisor of a
\emph{smooth} centre along which a sum of squares vanishes has order exactly $2$)
and $h_j+1=c$ for a codimension-$c$ smooth blow-up. Then a single binding divisor
of codimension $c$ contributes the ratio $c/2$. This is why, geometrically,
$\rlct=\tfrac12\min_t\codim S(t)$: each resolution branch's binding divisor has
codimension $\codim S(t)=\Mval(t)$. Keep this in view --- it is the antidote to
the per-layer recursion that the build drifted into.}

% =====================================================================
\section{The two theorems and Definition 3}
% =====================================================================

\subsection{Theorem 1: the reduced-rank-regression / three-layer ($L=2$) case}

This is Aoyagi--Watanabe 2005 (ref.\ [12]); the paper restates it as the $L=2$
benchmark that Theorem 2 generalises. Here $L=2$, two matrices $A^{(1)}A^{(2)}$,
widths $H^{(1)},H^{(2)},H^{(3)}$, reduced $M^{(s)}=H^{(s)}-r$.

\begin{definition}[the index set $\mathcal M$, $L=2$]\label{def:Mset-L2}
\[
 \mathcal M=
 \begin{cases}
   \{1,2,3\}, & \text{if } M^{(s)}<\tfrac12\sum_{l=1}^3 M^{(l)}\ \text{for all }s,\\
   \{1,2\}, & \text{if } M^{(3)}\ge M^{(1)}+M^{(2)},\\
   \{1,3\}, & \text{if } M^{(2)}\ge M^{(1)}+M^{(3)},\\
   \{2,3\}, & \text{if } M^{(1)}\ge M^{(2)}+M^{(3)};
 \end{cases}
\]
$\ell=\#\mathcal M-1$,\quad $\widetilde M=\tfrac1\ell\sum_{i\in\mathcal M}M^{(i)}$,\quad
$\Ms\in\mathbb N$ with $\Ms-1<\widetilde M\le \Ms$,\quad
$a=\sum_{i\in\mathcal M}M^{(i)}-(\Ms-1)\ell$.
\end{definition}

\begin{theorem}[Aoyagi Thm 1 = RRR; ref.\ {[12]}]\label{thm:rrr}
For the $L=2$ model, the learning coefficient of
$\norm{A^{(1)}A^{(2)}-\bar A^{(1)}\bar A^{(2)}}^2$ and its order are
\[
 \rlct=\frac{-r^2+r(H^{(1)}+H^{(3)})}{2}
       +\frac{a(\ell-a)}{4\ell}
       -\frac{\ell(\ell-1)}{4}\Big(\frac{\sum_{i\in\mathcal M}M^{(i)}}{\ell}\Big)^2
       +\frac12\!\!\sum_{\substack{i<j\\ i,j\in\mathcal M}}\!\! M^{(i)}M^{(j)},
 \qquad \rorder=a(\ell-a)+1.
\]
\end{theorem}

\fnote{\textbf{Deliverable 3.} We will (a) define the RRR model and \texttt{rlctRRR}
in Lean, (b) prove this value via the $L=2$ resolution (the machinery the build
already has for $L=2$, once re-anchored to the binding-divisor form), and
(c) once general-$L$ lands, expose Theorem~\ref{thm:rrr} as its $L=2$ instance.
The regular prefactor $\tfrac{-r^2+r(H^{(1)}+H^{(3)})}{2}$ is the Theorem~3
contribution; the rest is $\tfrac12\min_t\Mval(t)$ on the $L=2$ core.}

\subsection{Theorem 2: general $L$, and Definition 3}

\begin{definition}[Aoyagi Def.\ 3, general $L$]\label{def:Mset}
With $M^{(s)}=H^{(s)}-r$, $s=1,\dots,L+1$, choose
$\mathcal M=\{M^{(S_j)}:j=1,\dots,\ell+1\}\subseteq\{M^{(s)}\}$ such that
\[
 M^{(S_j)}<M^{(s)}\ \text{for } M^{(s)}\notin\mathcal M,\qquad
 \sum_{k=1}^{\ell+1}M^{(S_k)}>\ell M^{(s)}\ \text{for } M^{(s)}\in\mathcal M,\qquad
 \sum_{k=1}^{\ell+1}M^{(S_k)}\le(\ell-1)M^{(s)}\ \text{for } M^{(s)}\notin\mathcal M
 \ \ \text{(printed; see (T-D) --- this $\le$ is itself a suspected typo for $\ge$).}
\]
Let $\Ms$ with $\Ms-1<\tfrac1\ell\sum_{k=1}^{\ell+1}M^{(S_k)}\le\Ms$ and
$a=\sum_{k=1}^{\ell+1}M^{(S_k)}-(\Ms-1)\ell$.
\end{definition}

\tflag{\textbf{The underspecification (image-verified; see ledger (T-D)/(F-1)).}
Definition~3 \emph{prints} condition (iii) with ``$\le$'', and that print is itself a
\textbf{typo} (T-D): with ``$\le$'', non-members at $\ell=1$ require
$\sum_k M^{(S_k)}\le0$ --- impossible for positive widths --- so printed Def.~3
\emph{rejects exactly the size-2 sets Theorem~1 uses} (witness $M=[1,1,3]$: Thm~1 picks
$\{1,1\}$; printed Def.~3 rejects it), and is empty on the \emph{majority} of width
vectors ($105/216$ at $L{=}2$). The intended ``$\ge$'' reduces (max-$\ell$) to
Theorem~1's total $L=2$ selector but is then \emph{ambiguous in $\ell$}. So
\textbf{neither raw direction is a clean primitive.} \emph{Fix (brief standing
decision 3; verified sound, F-1):} take
$\rlct_{\mathrm{core}}:=\tfrac12\Mval_{\min}$, $\Mval_{\min}=\min_t\Mval(t)$ over the
\textbf{$t_L=0$} admissible-profile lattice (always non-empty) --- total,
permutation-invariant, $=\tfrac12(\sum_i q_i^2-\sum_k m_k^2)$, and matching the paper's
total $L=2$ RLCT on \emph{all} $216$ vectors (including the $105$ where printed Def.~3
is empty). Def.~3 is then a computed consequence where (under ``$\ge$'', max-$\ell$) it
selects a set --- never a primitive.}

\begin{theorem}[Aoyagi Thm 2, general $L$]\label{thm:main}
\[
 \rlct=\frac{-r^2+r(H^{(1)}+H^{(L+1)})}{2}
       +\frac{a(\ell-a)}{4\ell}
       -\frac{\ell(\ell-1)}{4}\Big(\frac{\sum_{j=1}^{\ell+1}M^{(S_j)}}{\ell}\Big)^2
       +\frac12\!\!\sum_{1\le i<j\le \ell+1}\!\! M^{(S_i)}M^{(S_j)},
\]
$\rorder=a(\ell-a)+1$.

The paper gives two further rewrites, all three verified algebraically equal by
exact arithmetic (\texttt{verify-arith-groundtruth.md}). Writing
$P:=\sum_{j=1}^{\ell+1}M^{(S_j)}=a+(\Ms-1)\ell$, so $P/\ell=\Ms+\tfrac{a-\ell}{\ell}$:
\begin{align*}
\rlct
&=\tfrac{-r^2+r(H^{(1)}+H^{(L+1)})}{2}
 +\tfrac{a(\ell-a)}{4\ell}-\tfrac{\ell(\ell-1)}{4}\big(\tfrac{P}{\ell}\big)^2
 +\tfrac12\!\!\sum_{i<j}\!M^{(S_i)}M^{(S_j)}\tag{A}\\
&=\tfrac{-r^2+r(H^{(1)}+H^{(L+1)})}{2}
 +\tfrac{a(\ell-a)}{4\ell}-\tfrac{\ell(\ell-1)}{4}\big(\Ms+\tfrac{a-\ell}{\ell}\big)^2
 +\tfrac12\!\!\sum_{i<j}\!M^{(S_i)}M^{(S_j)}\tag{B}\\
&=\tfrac{-r^2+r(H^{(1)}+H^{(L+1)})}{2}
 -\tfrac14(\ell-a-1)(\ell-a)-\tfrac{\ell(\ell-1)}{4}\big(\Ms^2+2\tfrac{a-\ell}{\ell}\Ms\big)
 +\tfrac12\!\!\sum_{i<j}\!M^{(S_i)}M^{(S_j)}.\tag{C}
\end{align*}
\tflag{(T-A) Transcription trap (\texttt{verify-arith}): form (B)'s middle term is
$\big(\Ms+\tfrac{a-\ell}{\ell}\big)^2$, \textbf{not} $\big(\tfrac{\Ms+a-\ell}{\ell}\big)^2$
--- $\Ms$ is not divided by $\ell$. The two are unequal; only the former equals (A).
Carry this when transcribing into Lean.}
\end{theorem}

\fnote{\textbf{Deliverable T (hero).} \texttt{aoyagi\_learning\_coefficient}.
We define \texttt{aoyagiLambda} $=\tfrac{-r^2+r(H^{(1)}+H^{(L+1)})}{2}+\rlct_{\mathrm{core}}$
with $\rlct_{\mathrm{core}}=\tfrac12\Mval_{\min}$ (the robust form), and prove it
equals Theorem~\ref{thm:main}'s expression where Def.~3 applies. The clean
internal target proven by the expedition's diagnostics is
$2\rlct_{\mathrm{core}}=\tfrac12\big(\sum_{i=1}^\ell q_i^2-\sum_{k=1}^{\ell+1}m_k^2\big)$
with $m_1\le\cdots\le m_{\ell+1}$ the $\ell+1$ smallest reduced widths and $q$ the
balanced $\ell$-part split of $P=\sum m_k$ --- this is \texttt{lambdaCore} and is
far more formalisable than the printed expression.}

\fnote{\textbf{Lean status (2026-07-21).} The wired, cite-free form is
\texttt{aoyagi\_learning\_coefficient\_via\_engine} (\texttt{DLN.Aoyagi.LearningCoefficient}):
$\mathrm{rlctGlobal}(\texttt{lossDLN}\,d\,0)=\texttt{cCodim}/2$ at the zero-product fibre,
for monotone positive widths, routed entirely through the engine (A$+$B$+$C$+$D) with
NEITHER \texttt{cited\_aoyagi\_lower\_ax} NOR \texttt{cited\_watanabe\_upper\_ax} on its
cone. Its axiom cone is clean-three $+$ \texttt{sorryAx} from exactly two leaves (the
min-over-charts CoV; the monument) --- when those close, the DLN lower-bound cite dies.
Named scope boundaries: $B=0$ (general $r$ = Skeleton lane, open); \texttt{Monotone d}
(non-monotone via banked permutation-invariance = a future leaf); the printed-Thm-2-form
equivalence (owed item O5).}

\begin{example}[equal widths]
If $M^{(1)}=\cdots=M^{(L+1)}$ then $\ell=L$, $\Ms-1<\tfrac{L+1}{L}M^{(1)}\le\Ms$,
$a=(L+1)M^{(1)}-(\Ms-1)L$, and
$\rlct=\tfrac{-r^2+2rH^{(1)}}{2}+\tfrac{a(L-a)}{4L}+\tfrac{L+1}{4L}(M^{(1)})^2$.
As $L\to\infty$, $\rlct\to\tfrac{-r^2+2rH^{(1)}}{2}+\tfrac{M^{(1)}}{4}+\tfrac14(M^{(1)})^2$:
\emph{bounded in depth} --- the paper's headline qualitative point.
\end{example}

\qmark{Cross-check the equal-width example against our ground truth:
$(2,2,2)$ i.e.\ $L=2,M^{(1)}=2,r=0$ should give $\rlct=3/2,\rorder=1$;
$(2,2,2,2)$ i.e.\ $L=3,M^{(1)}=2,r=0$ should give $\rlct=3/2,\rorder=3$. Verify the
closed form reproduces these (it must); this is the smallest sanity gate on the
arithmetic transcription.}

% =====================================================================
\section{Proof, part I: product reduction (Theorem 3)}
% =====================================================================

The proof has three movements: (I) peel the regular part of the product
(Theorem 3), reducing to the \emph{core} $\norm{\prod_s C^{(s)}}^2$ on reduced
widths; (II) reduce to the deepest singular point (Theorem 4); (III) resolve the
core by a recursive blow-up and read off $\rlct$ (\S\ref{sec:resolution}).

\subsection{Block elimination (Lemma 2)}

\begin{lemma}[Aoyagi Lemma 2]\label{lem:blockelim}
Let $A=\begin{psmallmatrix}A_1&A_2\\A_3&A_4\end{psmallmatrix}$ be $h_1\times h_2$ of
rank $r_1$, with $A_1$ a regular (invertible) $r\times r$ block, $r\le r_1$. Then
with
$Q_1=\begin{psmallmatrix}E_r&O\\-A_3A_1^{-1}&E_{h_1-r}\end{psmallmatrix}$,
$Q_2=\begin{psmallmatrix}E_r&-A_1^{-1}A_2\\O&E_{h_2-r}\end{psmallmatrix}$,
\[
 Q_1AQ_2=\begin{psmallmatrix}A_1&O\\O&C_4\end{psmallmatrix},\qquad
 C_4=-A_3A_1^{-1}A_2+A_4,\quad \rk C_4=r_1-r.
\]
Writing $F_3=-A_3A_1^{-1}$, $F_2=-A_1^{-1}A_2$, the transforms $Q_1,Q_2$ are
unipotent (regular).
\end{lemma}

\fnote{This is Rung L1 (\texttt{block\_elimination}). $C_4$ is the Schur
complement; the key for us is that $Q_1,Q_2$ are \emph{regular} (units of the
local ring), so by Lemma~\ref{lem:ideal} the RLCT is preserved and the variable
changes $A_2\mapsto F_2$, $A_3\mapsto F_3$, $A_4\mapsto C_4$ are a local analytic
isomorphism (unit Jacobian). The expedition's \texttt{RankNormalForm} /
\texttt{left/right\_normal\_form\_of\_*\_vanish} (\#159) is exactly this content.}

\subsection{Theorem 3: peeling the regular part}

\begin{theorem}[Aoyagi Thm 3]\label{thm:prodreduce}
Near a product of rank $r$, there are regular
$P_1=\begin{psmallmatrix}E_r&O\\F_3&E_{H^{(1)}-r}\end{psmallmatrix}$,
$P_2=\begin{psmallmatrix}E_r&F_2\\O&E_{H^{(L+1)}-r}\end{psmallmatrix}$ with
\[
 P_1\Big(\prod_{s=1}^{L}A^{(s)}\Big)P_2
   =\begin{psmallmatrix}C_1&O\\O&\prod_{s=1}^{L}C^{(s)}\end{psmallmatrix},
\]
$C_1$ regular $r\times r$, and $C^{(s)}$ of size $M^{(s)}\times M^{(s+1)}$.
\end{theorem}

\begin{proof}[Proof (reproduced, induction on the number of factors $S$)]
We can assume WLOG (Lemma~\ref{lem:ideal}, absorbing the regular part of the true
product) that $\prod_s\bar A^{(s)}=\begin{psmallmatrix}E_r&O\\O&O\end{psmallmatrix}$.
\emph{Base $S=1$.} Write $A^{(1)}=\begin{psmallmatrix}A^{(1)}_1&A^{(1)}_2\\A^{(1)}_3&A^{(1)}_4\end{psmallmatrix}$
with $A^{(1)}_1$ the $r\times r$ block; since the product has rank $r$ and is the
identity on the top-left, $\rk A^{(1)}_1=r$ (regular near the basepoint). Lemma~\ref{lem:blockelim}
gives $Q^{(1)}_1A^{(1)}Q^{(1)}_2=\begin{psmallmatrix}A^{(1)}_1&O\\O&C^{(1)}_4\end{psmallmatrix}$.
\emph{Step $S\to S+1$.} Suppose
$Q'_1\big(\prod_{s\le S}A^{(s)}\big)Q'_2=\begin{psmallmatrix}C'_1&O\\O&\prod_{s\le S}C^{(s)}\end{psmallmatrix}$
with $C'_1$ regular. Put $A'^{(S+1)}=Q_2'^{-1}A^{(S+1)}=\begin{psmallmatrix}A'_1&A'_2\\A'_3&A'_4\end{psmallmatrix}$.
Then
\[
 Q'_1\Big(\prod_{s\le S}A^{(s)}\Big)A^{(S+1)}
 =\begin{psmallmatrix}C'_1A'_1&C'_1A'_2\\ \prod_{s\le S}C^{(s)}A'_3 & \prod_{s\le S}C^{(s)}A'_4\end{psmallmatrix}.
\]
Because $Q'_1\begin{psmallmatrix}E_r&O\\O&O\end{psmallmatrix}=\begin{psmallmatrix}E_r&O\\F'_3&O\end{psmallmatrix}$,
the basepoint forces $C'_1A'_1$ regular. Left-multiply by
$Q''_1=\begin{psmallmatrix}E_r&O\\-\prod_{s\le S}C^{(s)}A'_3(C'_1A'_1)^{-1}&E\end{psmallmatrix}$
to clear the bottom-left, set $C^{(S+1)}=-A'_3A_1'^{-1}A'_2+A'_4$ (the next Schur
complement), and right-multiply by $Q''_2=\begin{psmallmatrix}E_r&-A_1'^{-1}A'_2\\O&E\end{psmallmatrix}$
to clear the top-right. This yields
$\begin{psmallmatrix}C'_1A'_1&O\\O&\prod_{s\le S+1}C^{(s)}\end{psmallmatrix}$, and the
accumulated $Q''_1Q'_1$, $Q'_2Q''_2$ are again unipotent. Induction completes at $S=L$.
\end{proof}

Subtracting the true product and using Lemma~\ref{lem:ideal} once more,
\[
 P_1\Big(\prod_s A^{(s)}-\begin{psmallmatrix}E_r&O\\O&O\end{psmallmatrix}\Big)P_2
 =\begin{psmallmatrix}C_1-E_r&-F_2\\-F_3&\prod_s C^{(s)}-F_3F_2\end{psmallmatrix},
\]
and the cross term $F_3F_2$ lies in the ideal of the regular generators
$\{C_1-E_r,F_2,F_3\}$, so it drops. Hence the \textbf{RLCT splits}:
\begin{equation}\label{eq:thm3-split}
 \rlct_{\bar w}\Big\langle\textstyle\prod_s A^{(s)}-\prod_s\bar A^{(s)}\Big\rangle
 =\frac{r^2+r(H^{(1)}+H^{(L+1)}-2r)}{2}
 +\rlct_{0}\Big\langle\textstyle\prod_s C^{(s)}\Big\rangle .
\end{equation}

\gnote{The regular term counts $\tfrac12\#\{C_1-E_r,F_2,F_3\}=\tfrac12\big(r^2+r(H^{(1)}-r)+r(H^{(L+1)}-r)\big)
=\tfrac12\big(r^2+r(H^{(1)}+H^{(L+1)})-2r^2\big)=\tfrac{-r^2+r(H^{(1)}+H^{(L+1)})}{2}$:
a regular (Morse) block of that many independent coordinates, RLCT $=\tfrac12\cdot\#$.
This is a genuine ``$\tfrac n2$'' Morse contribution --- and it occurs \emph{once},
at the top, not per layer. The per-layer ``$n/2$'' of the drifted build mis-placed
this single top-level regular term into every recursion node.}

\fnote{\eqref{eq:thm3-split} is Rung L2 (\texttt{product\_reduction}) composed with
the Fubini/regular-block split (the expedition's S1Fubini, adding
$n/2$, $n=r(H^{(1)}+H^{(L+1)})-r^2$ --- matches the prefactor exactly). After this
step, $B$ is gone and the entire remaining problem is the \emph{core}
$\rlct_0\langle\prod_s C^{(s)}\rangle$ at the origin on reduced widths $M^{(s)}$.
\textbf{All of \S\ref{sec:resolution} is about this core.}}

\fnote{\textbf{Lean status (strike wave 1, 2026-07-21): $r=0$ path LANDED; general
$r$ open, off the kill-path.} At the zero-product fibre ($B=0$, $r=0$) there is no
regular block and \eqref{eq:thm3-split} degenerates to the identity: the landed,
sorry-free \texttt{coreReduction} (\texttt{DLN.Aoyagi.LearningCoefficient}) proves
$\mathrm{rlctGlobal}(\texttt{lossDLN}\,d\,0)=\mathrm{rlctAt}(\sum(\texttt{coreGen})_i^2)\,0$
via the linear measure-preserving flatten (\texttt{exists\_flatten}, PROVEN), the
Frobenius identity \texttt{lossDLN\_zero\_eq\_coreLoss}, and the
$\mathbb R_{\ge0}^\infty\!\to\!\mathbb R$ carrier bridge. The \emph{general-$r$}
Theorem~3 (the gauge-slice normal form, Skeleton \texttt{product\_reduction} /
obligation \#44) remains a tracked-open sorry on the older $\mathrm{dlnLoss}$ carrier
--- consciously deprioritised: the charter corollary is the zero fibre.}

% =====================================================================
\section{Proof, part II + III: deepest point and the recursive blow-up}
\label{sec:resolution}
% =====================================================================

This section has three movements: Theorem 4 (deepest point, image-confirmed p.14),
the Cases 1\,\&\,2 recursive blow-up (\emph{the mountain}, to fill from images
p.15--21), and reading off the candidate $M_{s,k}$ (image-confirmed p.22).

\subsection{Theorem 4: reduction to the deepest singular point}\label{ssec:thm4}

\begin{theorem}[Aoyagi Thm 4; method of {[22]}, Aoyagi 2013]\label{thm:deepest}
Let $F_1,\dots,F_m$ be homogeneous in $w_1,\dots,w_j$ ($j\le d$), and let $\varphi$
be $C^\infty$ with
$\varphi(0,\dots,0,w^\ast_{j+1},\dots,w^\ast_d)\ge\varphi(w^\ast_1,\dots,w^\ast_d)$
and $\varphi w$ homogeneous in $w_1,\dots,w_j$ near
$(0,\dots,0,w^\ast_{j+1},\dots,w^\ast_d)$. Then
\[
 \rlct_{(0,\dots,0,\,w^\ast_{j+1},\dots,w^\ast_d)}(\ideal{F_1,\dots,F_m},\varphi)
 \ \le\
 \rlct_{(w^\ast_1,\dots,w^\ast_j,\,w^\ast_{j+1},\dots,w^\ast_d)}(\ideal{F_1,\dots,F_m},\varphi).
\]
\end{theorem}

\fnote{Rung D1 --- the \emph{only} place the brief permits leaning on a second paper
(Aoyagi 2013, [22]), and only ``to the extent D1 needs.'' Consequence used: the
entries of $\prod_s C^{(s)}$ are homogeneous in the $C$-variables, so the RLCT at
the origin ($c^{(s)}_{ij}=0$) is $\le$ the RLCT at any nearby basepoint. As the
\emph{global} learning coefficient is the infimum of the local RLCT over the fibre,
and the origin is the cone vertex of the homogeneous $\{\prod C=0\}$ (in the closure
of every stratum), the global coefficient is realised at $0$:
$\rlct_{\mathrm{global}}=\rlct_0\langle\prod_s C^{(s)}\rangle$. Hence set every
intermediate rank $r^{(s)}=r$ WLOG and resolve the core at $0$.}

\gnote{This is the rigorous form of ``the deepest point is the worst point'': a
homogeneity + lower-semicontinuity \emph{domination}, not a value computation ---
the \emph{same} content as the expedition's proven, value-free
\texttt{deepest\_le\_of\_homogeneous\_core}. So that lemma is correctly aimed and
survives the re-scope; it is what collapses the headline's $\inf_{w\in\mathrm{fibre}}$
to $\rlct_0$. (It does \emph{not}, by itself, give the value --- that is the
resolution below.)}

\fnote{\textbf{Lean status (strike wave 1, 2026-07-21): the corollary's instance is
PROVEN; the paper's general form is not formalised.} The landed
\texttt{Foundations.GlobalHomog} proves (i) invariance of the global and local RLCT
under a measure-preserving homeomorphism (no linearity, no homogeneity), and (ii)
\texttt{rlctGlobal\_eq\_rlctAt\_zero\_of\_homogeneous}: for continuous nonnegative $F$
homogeneous of degree $D$ in \emph{all} variables, $\mathrm{rlctGlobal}\,F=
\mathrm{rlctAt}\,F\,0$ --- riding the proven \texttt{deepest\_le\_of\_homogeneous\_core}
and closing the deepest-point step for the $B=0$ loss. Theorem~4 \emph{as printed}
(homogeneity in a sub-block $w_1,\dots,w_j$ only, arbitrary admissible $\varphi$) is
NOT formalised and nothing on the kill-path needs it; the paper's statement is
image-confirmed (p.14, this pass) but its proof (via [22]) is not reproduced --- our
build proves its own instance independently. Named scope boundary, not a hidden gap.}

\subsection{The recursive blow-up and the candidate thresholds}\label{ssec:blowup}

\emph{(Written from images p.14--p.22. This is the mountain and the exact locus of
the build's drift, so it is reproduced at the level of \emph{structure + exponent
bookkeeping}, which is what R1 must reproduce; the full chart-by-chart matrix
transforms are referenced to the pages.)}

\paragraph{The inductive invariant.} Write $M(S)=\min\{M^{(s)}:1\le s\le S\}$. The
recursion maintains, indexed by $(S,J)$ with $0\le S\le L$ and $0\le J\le M(S+1)$,
\[
 \Big\langle\prod_{s=1}^L C^{(s)}\Big\rangle
 =\Big\langle\operatorname{diag}(b_1,\dots,b_{M(S)})
   \begin{psmallmatrix}E_J&O\\O&D_J\end{psmallmatrix}\prod_{s=S+1}^L C^{(s)}\Big\rangle,
\]
where $D_J$ is the $(M(S)-J)\times(M^{(S+1)}-J)$ \emph{residual} block, and the
$b_i$ are monomials in the exceptional coordinates $u_{s,k}$ introduced so far:
$b_0=1$, $b_i=\big(\prod_{\widetilde t_{s,k}=i-1}u_{s,k}\big)\,b_{i-1}$. At the start
$S=J=0$, $D_0=\prod_s C^{(s)}$. The recursion has cleared $J$ unit pivots in layer
$S$; when $J$ reaches $M(S+1)=\min\{M(S),M^{(S+1)}\}$ the layer is done and $S$
increments. At $S=L+1$ the product is fully diagonal,
$\langle\prod_s C^{(s)}\rangle=\langle\operatorname{diag}(b_1,\dots,b_{M(L+1)})\rangle$,
and $\norm{\prod_s C^{(s)}}^2$ is (in these coordinates) $\sum_i b_i^2$, a sum of
squared monomials --- normal crossing.

\tflag{(T-F, image-verified 2026-07-21) \textbf{The paper's inductive statement claims
MORE than transcribed above --- and the extra claim is false.} The p.15 display (re-derived
at each case, pp.16/17/20) additionally asserts pairwise comparability of the carried
profiles: ``$T_{s,k}\le T_{s',k'}$ or $T_{s,k}\ge T_{s',k'}$'' for all pairs. This
totality is FALSE in general: at widths $(2,2,1,1)$ the profiles $(1,1,1)$ and $(2,1,0)$
are componentwise-incomparable (both $\Mval=1$; \texttt{g-monument-mval-instances.py}).
We deliberately do NOT carry it in the invariant here --- see the fidelity note at
\S\ref{sec:minimisation} for why nothing downstream needs it (the $\min$ needs no total
order; the $b$-chain, which does the work, is ordered by construction).}

\paragraph{The Jacobian.} Each exceptional divisor $u_{s,k}$ carries Jacobian power
$M_{s,k}-1$ (p.15): the change-of-variables determinant is
$\prod_{s,k}u_{s,k}^{\,M_{s,k}-1}\,du_{s,k}$ times the residual differentials. So in
the boxed S2 rule, $k_j\equiv1$ (each $b_i$ is a product of the $u$'s, the loss
$\sum b_i^2$ vanishes to order $2$ along each $u_{s,k}=0$) and $h_j=M_{s,k}-1$; the
ratio for divisor $u_{s,k}$ is $(h+1)/(2k)=M_{s,k}/2$.

\fnote{\textbf{Image-confirmed + the record correspondence (elder pass, 2026-07-21).}
The p.15 Jacobian-ledger display ($\prod u_{s,k}^{M_{s,k}-1}du_{s,k}$) and the $b_i$
recursion $b_0=1$, $b_i=\big(\prod_{\tilde t_{s,k}=i-1}u_{s,k}\big)b_{i-1}$ are now
IMAGE-VERIFIED against the page. Unrolling gives the thread-31 $\forall$-general closed
form $b_i=\prod_{\tilde t_{s,k}<i}u_{s,k}$ (structural induction, certificate \S b),
whence $b_1\mid b_2\mid\cdots\mid b_{M}$ \emph{by construction} --- the divisibility
chain that is the Lean record's \texttt{hchain} and the hypothesis under which Object~C's
boxed rule is a theorem. The Lean \texttt{Chart} record separates the two exponent
ledgers the paper interleaves: \texttt{bexp} (the loss-side monomial exponents, here
$k_j\equiv1$ on binding axes $=$ \texttt{hunit\_mult}) vs \texttt{jac} (the
Jacobian-side $h_j=M_{s,k}-1$, certified against the actual $|\!\det Dg|$ by
\texttt{hjac}); conflating them was a v2-blueprint defect, killed by counterexample.}

\paragraph{The step (Cases 1 \& 2).} At $(S,J)$, look at the equal run of the
$b$-sequence above $J$:
\begin{itemize}
\item \textbf{Case 1} (p.15--18): $b_{J+1}=\dots=b_{J+J_1}\ne b_{J+J_1+1}$, a
\emph{partial} equal block of length $J_1<M(S)-J$. Blow up
$\{d_{ij}=0\ (J{<}i\le J{+}J_1,\ J{<}j\le M^{(S+1)}),\ u_{s,k}=0\}$. The chart splits:
\emph{Case 1(1)} (the $d$-block $=u_{s,k}\cdot d'$) sets $\widetilde t_{s,k}{=}J$ and
\textbf{adds} $M'_{s,k}=M_{s,k}+J_1(M^{(S+1)}-J)$ to that divisor's exponent
(decrementing the count of $u$'s with $\widetilde t=J{+}J_1$ by one --- an inner
recursion); \emph{Case 1(2)} (first row normalised, $u_{s,k}=u_{S,J+1}u'_{s,k}$)
introduces $u_{S,J+1}$, and regular transforms $Q,P$ reduce $D_J''$ to
$\begin{psmallmatrix}1&O\\O&D_{J+1}\end{psmallmatrix}$, advancing $J$ (or $S$).
\item \textbf{Case 2} (p.19--22): $b_{J+1}=\dots=b_{M(S)}$, the \emph{full} remaining
block. Blow up $\{d_{ij}=0\ (J{<}i\le M(S),\ J{<}j\le M^{(S+1)})\}$. The chart
($d$-block $=u_{S,J+1}\cdot d'$, top-left $1$) sets $t^{(i)}_{S,J+1}=M^{(i+1)}$
($i{<}S$), $t^{(S)}_{S,J+1}=\dots=t^{(L)}_{S,J+1}=J$, and the divisor exponent
\[
 M'_{S,J+1}=(M(S)-J)(M^{(S+1)}-J)
\]
--- the codimension of the residual block being blown up. Regular $Q,P$ then reduce
$D_J''\to\begin{psmallmatrix}1&O\\O&D_{J+1}\end{psmallmatrix}$, advancing $J$ (or $S$).
\end{itemize}

\tflag{(T-C, PINNED --- elder image pass 2026-07-21) The Case-2 exponent
$M'_{S,J+1}=(M(S)-J)(M^{(S+1)}-J)$ is confirmed against the p.20 image, and the
Case-1 exponents $M'_{s,k}=M_{s,k}+J_1(M^{(S+1)}-J)$ (Case 1(1), p.16) and
$M'_{S,J+1}=M_{s,k}+J_1(M^{(S+1)}-J)$ (Case 1(2), p.17) likewise. Non-bindingness of
the Case-2 divisor at the flagged instance is battery-verified
(\texttt{g-monument-mval-instances.py}, ``(2,2,3,2) defect non-binding'' PASS).}

\tflag{(T-E, image-CONFIRMED 2026-07-21) \textbf{The Case-2 head-reset label is RAW-width
--- and asymmetric with Case 1(2).} p.20 prints $t^{(i)}_{S,J+1}=M^{(i+1)}$ for $i<S$
(the RAW width), while Case 1(2) at p.17 prints $t^{(i)}_{S,J+1}=t^{(i)}_{s,k}$
(INHERITED from the factored divisor). At non-monotone widths the raw label conflicts
with the paper's own p.22 formula --- the defect ledgered in
\texttt{verify-case2-rawwidth-defect.md}, previously text-only, now confirmed on the
page image. RESOLUTION (thread-31, structural): the ideal-route Lean record carries
\emph{no label field at all} (only \texttt{bexp} + \texttt{jac}); the exponent
accumulation is governed by the RUNNING-MIN $M(S)$, so the defect has no formal
representation to corrupt; at the witness $(2,2,3,2)$ the defect divisor is
non-binding. The p.20-image residual of the compass Case-2 ledger entry is CLOSED.}

\gnote{\textbf{The coupled mechanism is VERIFIED, and the single chain SURVIVES
(threads 27/28, exact algebra + Gr\"obner + decorrelated Codex).} The genuinely
coupled corank-$(2,2)$ instance $(3,3,4)$, $t=(1,0)$, $\Mval=8$ (coupled-only
minimiser --- no clean peel reaches it): the peel$+$radial$+$join mechanism produces
exactly the $\operatorname{diag}(b)$ of this recursion with $\rlct=4=\tfrac12\Mval$
(\texttt{g-coupled-334-diagb.py}). Two independent shared deep factors do NOT break
principality: at $(4,4,4)$, $t=(2,0)$ (two corank-2 blocks) and the $L=4$ separated
instance $(3,3,3,2,2)$, $t=(2,2,1,0)$, the terminal chart's ideal is again a single
principal chain (corner-join restores principality; \emph{intermediate} charts can be
non-principal --- the chain is a TERMINAL-chart invariant), verified by Gr\"obner
ideal-equality, not coefficient-matching. So the paper's single-chain inductive
invariant survives every kill-instance the expedition could construct; the flagged
open end is documented in the thread-28 certificate (deeper mixed instances), folded
into the monument seat's obligations.}

\subsection{The candidate thresholds and the value}\label{ssec:candidates}

Carrying the exponents through to $S=L+1$, each terminal divisor (those with
$\widetilde t_{s,k}=0$) has accumulated exponent (p.22)
\[
 M_{s,k}=(M^{(1)}-t^{(1)}_{s,k})(M^{(2)}-t^{(1)}_{s,k})
        +\sum_{j=2}^{L}(t^{(j-1)}_{s,k}-t^{(j)}_{s,k})(M^{(j+1)}-t^{(j)}_{s,k}),
\]
indexed by the rank-profile $t_{s,k}=(t^{(1)},\dots,t^{(L)})$ of its branch, and
\[
 \boxed{\ \rlct_{\mathrm{core}}=\tfrac12\,\min\{M_{s,k}:\widetilde t_{s,k}=0\}.\ }
\]

\gnote{$M_{s,k}=\Mval(t_{s,k})=\codim S(t_{s,k})$ \emph{exactly} (thread-03, proven
general-$L$). So $\rlct_{\mathrm{core}}=\tfrac12\min_t\codim S(t)=\tfrac12\Mval_{\min}$
--- the value the red-team certified (\S\ref{ssec:thm4} ledger F-1), and a
\emph{direct min over per-branch full-codimension divisor exponents}. \textbf{There
is no $\min\{mk/2,\,n/2+\text{child}\}$ per-layer RLCT recursion} --- that was the
build's drift.}

\fnote{\textbf{Lean status (2026-07-21): the read-off layer is built; the two remaining
sorries are exactly this display's two halves.} The $\tilde t_{s,k}=0$ restriction is
\texttt{bindingAxes} (positive exponents of the dominant $b_{k_0}$); the per-chart min is
\texttt{Chart.chartMin}; the atlas min is \texttt{Resolution.divisorMin}; and Object D's
bridge $\texttt{divisorMin}=\texttt{qipMin}=\texttt{cCodim}$ is PROVED
(\texttt{Core.Aoyagi.Engine.divisorMin\_eq\_cCodim}, min-attainment form). The whole
engine's remaining sorried surface is: (1) the min-over-charts change-of-variables
\texttt{rlctAt\_sumSqFam\_eq\_iInf\_charts} (the boxed equality itself), and (2)
\texttt{exists\_coreResolution} (that THIS recursion, at coupled corank $\ge2$, inhabits
the certified atlas record for the concrete flattened core --- the monument; its
combinatorial half, $\Mval$-attainment across leaves, is salvaged kernel-checked from the
retired Engine's tree under elder conditions C1--C2; the geometric fields are the true
remaining build).}

\paragraph{R1 design verdict (\#18; \texttt{verify-r1-shortcut.md}, decorrelated +
Codex, exact).} The design-\S8 \emph{one-shot} shortcut (iterated block-elim to the
codim-$\Mval(t)$ residual $+$ \emph{one} blow-up $\Rightarrow F=u^2\cdot\text{unit}$)
is \textbf{REFUTED for $L\ge3$} --- the same failure class as the per-node recursion.
Witness ($L{=}3$, $(2,2,2,2)$, $t{=}(1,0,0)$, $\Mval{=}3$): the $L{=}2$ incidence chart
$A=\alpha\!\begin{psmallmatrix}1&a\\b&ab+\delta\end{psmallmatrix}$,
$B=\begin{psmallmatrix}u-ar&v-as\\r&s\end{psmallmatrix}$, then one blow-up of
$\{\delta{=}u{=}v{=}0\}$, gives $F=\alpha^2\rho^2\,U$ with
$U=\norm{\begin{psmallmatrix}\xi&\eta\\r&s\end{psmallmatrix}C^{(3)}}^2$ \textbf{vanishing
to order $4$} --- NOT a unit. (At $L{=}2$ the residual is a nondegenerate rank-$4$ Morse
quadratic; that is precisely why $L{=}2$ works and only $L{=}2$.)

\textbf{Corrected R1 route --- a depth-recursion} (Aoyagi's actual layer-peeling,
streamlined: no $b_i$-product / Case-1(1) exponent-merging bookkeeping). One
(incidence $+$ one blow-up) peels \emph{exactly one layer}:
$F=\alpha^2\rho^2\cdot\norm{X\,C^{(3)}\cdots C^{(L)}}^2$, a \emph{fresh depth-$(L{-}1)$
core}. Verified exact at $L{=}3,4$, rank-$1$ peel, incl.\ non-square inner widths
$(2,2,3,2)$. The value $\rlct_{\mathrm{core}}=\tfrac12\min_t\Mval(t)$ is \textbf{robust};
the cover lower bound (no branch ratio $<\tfrac12\min_t\Mval$) is structural/inductive
($\alpha$-divisor $\ge\tfrac12\Mval(0)$, $\rho$-divisor $=\tfrac12\Mval(\text{branch})$,
recursive divisors $\ge\tfrac12\min$ by induction) --- no full atlas needed.
\gnote{This IS Aoyagi's depth recursion (peel a layer $\to$ lower-depth core), with
\emph{exact} per-step geometry --- \emph{not} the per-node Morse recursion (whose
$n{=}M(\text{last})$ overshot) and \emph{not} the \S8 one-shot. R1 builds THIS.}
\paragraph{Partial-rank leg --- CLOSED (\texttt{verify-r1-shortcut}, exact + Codex).}
The clean disjoint ``divisor $+$ \emph{independent} fresh core'' is a \textbf{rank-1
($t_1{=}1$) phenomenon only}. For partial rank ($t_1\ge2$; witness $(3,3,2,2)$ binding
branch $t{=}(2,1,0)$, $\Mval{=}4$, $\rlct{=}2$): block-eliminating
$C^{(1)}\to\operatorname{diag}(E_2,\delta)$ gives
$F_{\mathrm{core}}=\norm{T\,C^{(3)}}^2+\delta^2\norm{R\,C^{(3)}}^2$, and the two terms
\emph{share} $C^{(3)}$ --- after resolving $\norm{TC^{(3)}}^2$ the
$\delta^2\norm{RC^{(3)}}^2$ term is \emph{not} divisible by the exceptional coordinate.
The peel leaves $\operatorname{diag}(E_2,\delta)\cdot(\text{free})$, \textbf{exactly
Aoyagi's $\operatorname{diag}(b)$ invariant}: the recursion continues on a
\emph{constrained/coupled} core, it does not factor. (Codex: the shared-$C^{(3)}$
coupling \emph{raises} the residual threshold from $3/2$ to exactly $2$ --- load-bearing.)
\gnote{\textbf{The value $\rlct_{\mathrm{core}}=\tfrac12\min_t\Mval(t)$ is robust
general-$L$, all branch types} (ideal-preservation: block-elim is a unit transform
$\Rightarrow\rlct=\tfrac12\Mval(\text{branch})$; $+$ the structural cover lower bound).
Confirmed $=2$ on $(3,3,2,2)$. \textbf{But the clean-disjoint MECHANISM is
$t_1{=}1$-only}: R1's \emph{proof} must use Aoyagi's coupled $\operatorname{diag}(b)$
recursion (option (ii)) for general branches; the clean disjoint Fubini recursion
(option (i)) covers only $t_1{=}1$ minimizers. \textbf{R1 design: build option (ii)}
(Aoyagi's actual coupled recursion); the headline value is safe regardless.}
\paragraph{The partial-rank dichotomy --- the final characterisation
(\texttt{verify-r1-shortcut}, both named cases run exactly).} By the layer-1 residual codimension
$c_1=(M^{(1)}-t_1)(M^{(2)}-t_1)$:
\begin{itemize}
\item \textbf{$c_1=0$ (rank full in the smaller dimension): clean, no layer-1 divisor.} Either
$t_1=M^{(2)}$ (full column rank: $C^{(1)\top}C^{(1)}$ is a unit, $\langle C^{(1)}\cdots\rangle=\langle
C^{(2)}\cdots\rangle$) or $t_1=M^{(1)}$ (full row rank: $C^{(2)}\mapsto C^{(1)}C^{(2)}$ surjective,
Jacobian full) --- both reduce to a \emph{fresh lower-depth core} by a unit transform. Witnesses
$(3,2,2,2)\,t_1{=}2$ and $(2,3,2,2)\,t{=}(2,1,0)$ both land here ($\rlct=3/2=\tfrac12\min\Mval$).
\item \textbf{$0<t_1<\min(M^{(1)},M^{(2)})$ (genuine partial drop): coupled.} The $\operatorname{diag}(b)$
invariant; the disjoint factorisation fails (witness $(3,3,2,2)$, $c_1{=}1$, value $=2$).
\end{itemize}
\gnote{\textbf{Value certified general-$L$, all branch types} (rank-1, $c_1{=}0$, $c_1{>}0$): every case
evaluates to exactly $\tfrac12\Mval(\text{branch})$ by ideal-preservation, so
$\rlct_{\mathrm{core}}=\tfrac12\min_t\Mval(t)$ --- safe to build the headline on. The clean disjoint
\emph{mechanism} holds for $c_1{=}0$ and rank-1 peels; $c_1{>}0$ needs the coupled recursion.}
\paragraph{The light-vs-coupled fork --- RESOLVED: R1 must carry the coupled $\operatorname{diag}(b)$
support (\texttt{verify-r1-light-recursion.md}, \#18 cont'd; Newton-LP $+$ decorrelated Codex, exact to
the rational).} A \emph{threshold-only} invariant (residual widths $+$ a per-output-row weight
\emph{multiplicity}, no symbolic matrix support) is \textbf{provably insufficient}. It is faithful only
when every binding-branch peel has $\operatorname{corank}\le1$ (the residual block $\Delta$ is a scalar,
its monomial support \emph{forced} by widths$+$branch$+$weight) --- covering the rank-$1$ chains,
$(3,2,2,2)$, and even $(3,3,2,2)\,t{=}(2,1,0)$ ($\Delta=\delta$ scalar). It \textbf{BREAKS at
$\operatorname{corank}\ge2$} (genuine $\ge2\times2$ $\Delta$-block): a per-row multiplicity cannot encode
\emph{which} divisor variables are SHARED, and sharing changes the Newton polytope hence the RLCT. Exact
obstruction (two independent computations agree): $\ideal{\delta x,\delta y}$ has $\rlct=\tfrac12$ vs
$\ideal{\delta_1 x,\delta_2 y}$ has $\rlct=1$ --- \emph{identical} light data, different value.
\textbf{Correction (\texttt{verify-r1-diagb-4422}; exact $\times3$ methods $+$ decorrelated Codex):
$(4,4,2,2)$ is NOT the binding witness.} Its $\rlct$ is $2$, not $\tfrac72$ --- the \emph{clean} branch
$t{=}(4,2,0)$ ($\Mval{=}4$) binds via one radial $C^{(3)}$ blow-up; the corank-$2$ branch $t{=}(2,1,0)$
($\Mval{=}7$) is \emph{non-binding} (the RLCT is the min over \emph{all} branches). So that branch does
not force the coupled structure (and threshold-only there gives the \emph{correct} $2$). The genuinely
coupled-binding witness is \textbf{$(3,3,4)$ [an $L{=}2$ RRR core]}: $\Mval(t_1)=(3{-}t_1)^2+4t_1=9,8,9,12$,
so the minimiser is $t{=}(1,0)$, $\Mval{=}8$, $\rlct{=}4$, layer-1 corank $(2,2)$ --- and the clean
branches $t_1{\in}\{0,3\}$ give $9,12>8$, so \emph{no} $\operatorname{corank}\le1$ peel reaches the
minimiser. \textbf{This is now CERTIFIED} (\texttt{verify-r1-diagb-334.md}, PROVED $+$ decorrelated Codex
$+$ Aoyagi--Watanabe 2005 RRR anchor): the explicit coupled $\operatorname{diag}(b)$ resolution gives
$\rlct{=}4{=}\tfrac12\Mval$, while a threshold-only (per-row multiplicity) recursion gives $3$ --- so the
obstruction \emph{binds} at a genuinely-binding DLN branch, and the symbolic $\operatorname{diag}(b)$
support is NECESSARY (the test $(4,4,2,2)$ failed to provide, its binder being clean).
\gnote{\textbf{R1 Lean shape --- DECIDED: option (ii), the coupled $\operatorname{diag}(b)$ recursion}
(Aoyagi's actual structure). Minimal sufficient invariant $=$ the per-generator \emph{symbolic divisor
support $+$ sharing relations}, up to unit equivalence (numeric coefficients / analytic unit changes need
not be carried). The ``hybrid'' option (i) --- clean for $\operatorname{corank}\le1$, corank-$\ge2$
branches ``by the value $=\tfrac12\codim$'' --- is \textbf{not one-citation-viable}: it would import the
forbidden $\rlct=\tfrac12\codim$ cite for exactly the hard branches (brief: S2 only). So the coupled
recursion is the \emph{only} sound general-$L$ route under our constraint. This \textbf{vindicates} the
build's \texttt{GeneralR1Recursion}/routeStep dispatcher (whose SOUNDNESS NOTE already records that the
dimension-conserving \emph{clean} recursion is false --- telescopes to ambient$/2{=}4\ne\tfrac32$ for
$(2,2,2)$): the coupled recursion is necessary, not over-engineering. (The DECISION rests on the abstract
$\ideal{\delta x,\delta y}$-vs-$\ideal{\delta_1 x,\delta_2 y}$ obstruction $+$ clean-reachability failing
at a genuine corank-$2$ minimiser --- confirmed for $(3,3,4)$, \emph{not} $(4,4,2,2)$ per the correction
above.) \emph{The Lean build's spec is now CERTIFIED} (\texttt{verify-r1-diagb-334.md} §3): the
$(3,3,4)\,t{=}(1,0)\to4$ binding resolution exhibits the exact symbolic support the routeStep datum must
carry --- \texttt{support : Gen}$\to$\texttt{Finset DivVar} (which exceptional $u$-vars divide which
generator) with the sharing identity load-bearing (it raises the $\Delta$-block contribution from $1$ to
$2$). \emph{Optional refinement (not blocking):} an $L\ge3$ binding witness exercising a \emph{shared
deep factor} $C^{(s)}$ rather than the $\Delta$-internal shear --- the support field already accommodates
it (the non-binding $(4,4,2,2)\,t{=}(2,1,0)$ demo exhibited the shared-$C^{(3)}$-via-$u$ structure).}

\gnote{\textbf{The earlier ``hard design questions'' are now settled} (\#18, decorrelated $+$ Codex):
(1) the \S8 single-binding-divisor one-shot does \emph{not} reproduce $\min M_{s,k}$ for $L\ge3$ (refuted,
the $(2,2,2,2)$ order-4 witness above); the depth-recursion / coupled recursion does. (2) The Case-2
printed exponent $(M(S)-J)(M^{(S+1)}-J)$ is the codimension of the blown-up block (image p.20, (T-C)) and
is confirmed \emph{non-binding} ($2\rlct\le M^{(i)}M^{(j)}$, never attains the min). No open validation leg
remains; R1 can build.}

% =====================================================================
\section{Proof, part IV: the minimisation and the closed form (Lemmas 3--5)}
\label{sec:minimisation}
% =====================================================================

Part~II gave, per terminal branch (rank-profile $t$), the divisor exponent $M_{s,k}=\Mval(t)$ and
$\rlct_{\mathrm{core}}=\tfrac12\min_t M_{s,k}$. Turning this into the closed form of
Theorem~\ref{thm:main} is two nested minimisations: an \emph{inner} balance within a fixed selected set
(Lemma~3), and the \emph{outer} selection of the set itself (Definition~3 / the $\ell+1$ smallest widths).

\emph{Fidelity note (documented paper defect, 2026-07-21; wording sharpened after the elder's page-image
pass): the rank-profiles carried by the recursion are NOT totally comparable, contrary to the paper's
EXPLICIT claim --- the totality ``$T_{s,k}\le T_{s',k'}$ or $T_{s,k}\ge T_{s',k'}$'' is a displayed part
of the inductive statement (p.15) and is re-derived in each case (pp.16/17/20), not merely an implicit
reading. It is false: at widths $(2,2,1,1)$ the profiles $(1,1,1)$ and $(2,1,0)$ are
componentwise-incomparable with equal $\Mval=1$ (witness \texttt{g-monument-mval-instances.py}; found by
a decorrelated consult, verified by three independent runs). No conclusion above depends on totality: the
minimisation is over the \emph{set} of admissible profiles (a $\min$ needs no total order on its index
set), and the divisibility chain $b_1\mid\cdots\mid b_M$ --- Aoyagi's own invariant, which does the real
work --- is ordered by construction and survives. Only an argument leaning on profile comparability
itself would be unsound; the formalisation's record carries no profile field at all (thread-31
certificate \S d). Ledger item (T-F); claim-site cross-reference at the inductive statement,
\S\ref{ssec:blowup}.}

\subsection{Lemma 3: the within-set balance (image-confirmed p.24)}\label{ssec:lemma3}

Fix the selected set $\mathcal M=\{M^{(S_1)},\dots,M^{(S_{\ell+1})}\}$, $P:=\sum_k M^{(S_k)}$, and recall
$\Ms-1<P/\ell\le\Ms$, $a=P-(\Ms-1)\ell$. The terminal branches that can bind the minimum differ only in
how Aoyagi's local-coordinate construction distributes the split; parametrised by an integer
$b\in\{0,\dots,\ell-1\}$ (Lemma~3's hypothesis is $0\le a,b\le\ell-1$, so $\ell-a\ge1$), the bound is the
one-variable quadratic (p.24)
\[
 A(b)=b(\ell-a)^2+(\ell-1-b)a^2+\big(b(\ell-a)-(\ell-1-b)a\big)^2
      =\ell^2 b^2+\ell^2(1-2a)b+a^2\ell(\ell-1),
\]
\[
 \frac{\partial A}{\partial b}=\ell^2(1+2b-2a),\qquad
 \min_{b\in\{0,\dots,\ell-1\}}A(b)=A(a-1)=A(a)=a\,\ell(\ell-a).
\]
The real minimiser is $b=a-\tfrac12$, so the integer neighbours $b=a-1,a$ tie (this tie is exactly what
$\rorder$ counts, \S\ref{ssec:order}). Dividing by $\ell^2$ and scaling by $\tfrac14$ recovers the first
variable term of $\rlct$:
\[
 \tfrac14\cdot\frac{\min_b A(b)}{\ell^2}=\frac{a(\ell-a)}{4\ell}.
\]
\fnote{The only genuinely-analytic minimisation in the value layer, and it is a one-variable integer
quadratic --- trivial in Lean once stated. Exact arithmetic $+$ decorrelated Codex confirmed the print
$A(a)=a\ell(\ell-a)$ is \emph{correct} (\texttt{verify-arith-groundtruth.md} Check~5); an earlier draft's
``typo'' (T-B) was retracted after reading p.24. Domain subtlety only: $A(a-1)$ needs $a\ge1$; at $a=0$
the lone admissible $b=0$ still gives $0=a\ell(\ell-a)|_{a=0}$.}

\fnote{\textbf{Lean status (2026-07-21).} In the engine, Lemma~3's role is carried by the
\emph{min-attainment} obligations of the resolution record (elder v2-defect-4 form): every binding
divisor's integer exponent is $\ge\texttt{qipMin}$ and some chart's divisor attains it
(\texttt{hlb}/\texttt{hattain} in \texttt{exists\_coreResolution}); the banked QIP layer
(\texttt{cCodim\_eq\_qipMin}, \texttt{minAdm\_eq\_cCodim}) then closes the loop to the codimension. The
salvaged Engine combinatorics (\texttt{minAdm\_le\_terminalExponents}, \texttt{o5\_core\_realized},
kernel-checked at all widths) is the adapter-side discharge of exactly these two obligations.}

\subsection{The set selection and the closed form}

The outer minimisation chooses $\mathcal M$ (hence $\ell$). Geometrically (\S\ref{ssec:candidates}) it is
$\min_t\Mval(t)$ over the admissible rank-profile lattice; combinatorially (Def.~3, modulo the (T-D)/(F-1)
underspecification of \S\ref{def:Mset}) the binding $\mathcal M$ is the set of the $\ell+1$ \emph{smallest}
reduced widths at the balancing $\ell$. Substituting the inner value gives Theorem~\ref{thm:main}:
\[
 \rlct=\underbrace{\frac{-r^2+r(H^{(1)}+H^{(L+1)})}{2}}_{\text{Thm 3 regular block}}
       +\underbrace{\frac{a(\ell-a)}{4\ell}-\frac{\ell(\ell-1)}{4}\Big(\frac P\ell\Big)^2
        +\frac12\!\!\sum_{i<j}\!M^{(S_i)}M^{(S_j)}}_{=\ \tfrac12\min_t\Mval(t)\ =\ \rlct_{\mathrm{core}}}.
\]

\subsection{The clean form (the formalisation target)}\label{ssec:clean}

The expedition repackages $\rlct_{\mathrm{core}}$ into a form with no $\Ms$, no $\lceil\cdot\rceil$, and
small terms. Let $m_1\le\cdots\le m_{\ell+1}$ be the $\ell+1$ smallest reduced widths, $P=\sum_k m_k$, and
$q=(q_1,\dots,q_\ell)$ the balanced $\ell$-part split of $P$ (each $q_i\in\{\lfloor P/\ell\rfloor,\lceil
P/\ell\rceil\}$, $\sum_iq_i=P$). Then
\[
 \boxed{\ 2\,\rlct_{\mathrm{core}}=\tfrac12\Big(\sum_{i=1}^{\ell}q_i^2-\sum_{k=1}^{\ell+1}m_k^2\Big).\ }
\]
This equals the bracketed closed form above for every width vector (exact enumeration $L\le4$ $+$ sampled
$L=5,6$; \texttt{verify-def3-underspec.md} Check~4), and is \texttt{lambdaCore} in Lean.

\fnote{\textbf{Lean status + an honestly-owed equivalence (2026-07-21).} The engine's landed headline
speaks in $\texttt{cCodim}$ (the geometric/QIP codimension): $\mathrm{rlctGlobal}=\texttt{cCodim}/2$
modulo the two named sorries. The chain $\texttt{cCodim}=\texttt{qipMin}=\texttt{minAdm}$ is PROVEN.
What is NOT proven $\forall(L,M)$ is the identification of $\texttt{cCodim}$ with the \emph{printed}
Theorem-2 expression (equivalently with this clean $q^2-m^2$ form): that equivalence is
exact-enumeration-verified ($L\le6$, bounded widths) only --- the (F-1) \emph{still-owed theorem}. It is
a purely combinatorial integer-program identity, consciously deferred: the destination names $C=$ the
codimension, which IS \texttt{cCodim}; the printed form is downstream exposition. Ranked in
\texttt{verify-owed-math-audit.md} (item O5).}
\fnote{\textbf{Why this is the Lean primitive, not Def.~3.} (i) \emph{Total}: the $\ell+1$ smallest widths
and a balanced split exist for \emph{every} width vector, where Def.~3's inequalities are empty on the
majority ($105/216$ at $L{=}2$; \S\ref{def:Mset}). (ii) \emph{Permutation-invariant}: $\Mval(t)$ singles
out $M^{(1)},M^{(2)}$ yet $\Mval_{\min}$ is identical on every ordering of the width multiset (0 violations,
$L\le4$) --- licenses sorting in \texttt{aoyagiLambda}. (iii) \emph{Small terms}: $\le 2(\ell+1)$ squares,
keeping \texttt{decide}/\texttt{norm\_num} cheap. Define \texttt{aoyagiLambda} $=$ (Thm-3 prefactor) $+
\tfrac12\Mval_{\min}$; prove $=$ the clean form (\texttt{lambdaCore\_eq\_clean}) and $=$ the printed
expression where Def.~3 applies. \emph{Do not transcribe Def.~3's inequalities as the definition.}}

\subsection{Lemmas 4--5: the order $\rorder=a(\ell-a)+1$ (images p.25--26)}\label{ssec:order}

The order $\rorder$ --- the multiplicity of the largest pole of the zeta function
$\int|F|^{-2c}\varphi\,dw$, geometrically the number of top-dimensional components of the exceptional fibre
over the binding stratum --- is counted combinatorially (p.25--26). Write $M:=\Ms$. The minimising
terminal branch-vectors are pinned between \emph{two} two-piece partial-sum envelopes:
\begin{align*}
 \widetilde H_j &= \textstyle\sum_{l=1}^{j+1}M^{(S_l)}-jM &&(1\le j\le a), &
 \widetilde H_j &= \textstyle\sum_{l=1}^{j+1}M^{(S_l)}-aM-(j-a)(M-1) &&(a< j\le\ell),\\
 \widetilde H'_j &= \textstyle\sum_{l=1}^{j+1}M^{(S_l)}-j(M-1) &&(1\le j\le\ell-a), &
 \widetilde H'_j &= \textstyle\sum_{l=1}^{j+1}M^{(S_l)}-(\ell-a)(M-1)-(j-\ell+a)M &&(\ell-a< j\le\ell),
\end{align*}
with only the \emph{terminal} equality $\widetilde H_\ell=\widetilde H'_\ell=0$ (\emph{not}
$\widetilde H_1=0$ in general --- e.g.\ $(2,2,2)$ gives $\widetilde H_1=1$). Lemma~4 characterises the
binding branch-vectors $T_{s,k}$ by \textbf{two} conditions: the envelope
$\widetilde T\le T_{s,k}\le\widetilde T'$ \emph{and} the local-increment condition
$H_{j-1}-H_j+M^{(S_{j+1})}\in\{M-1,M\}$ (the second forces the $a$ steps of size $M$ and the $\ell-a$
steps of size $M-1$). Lemma~5 then bounds the count \emph{above} by the union
$\bigcup_{j=1}^{\ell-1}\{H:\widetilde H_j\le H\le\widetilde H'_j\}$ together with the Case-1(2)
$J$-increment, giving $\rorder\le a(\ell-a)+1$, and an explicit local-coordinate construction attains it
(a two-sided argument, \emph{not} a single interval count). Hence
\[
 \rorder=a(\ell-a)+1.
\]
\gnote{This is the \emph{combinatorial} order --- a count of binding branches (for $\ell=1$,
$a=1$, it is $1$; the unique minimiser). Equating it to the \emph{analytic} pole multiplicity is the step
that needs meromorphic continuation, which Mathlib lacks.}

\fnote{\textbf{Lean status (2026-07-21): deferred per charter \S1-E, seams named.} Built:
\texttt{boxedOrder} (the per-chart count of min-attaining binding divisors) with
\texttt{one\_le\_boxedOrder} PROVEN, and the closed-form count \texttt{aoyagiPoleOrder}
$=a(\ell-a)+1$ + the count-vs-order distinction (\texttt{ThetaOrderDistinction}: $\rorder\neq$
the fibre component count \texttt{numTop}). NOT built, honestly named: (i) the combinatorial identity
boxedOrder-aggregate $=a(\ell-a)+1$ (needs the Lemma-3 $(\ell,a)$-balancing tie to the record); (ii)
the analytic identification with the zeta pole multiplicity (no pole-order operation exists to state it
against). E stays a reachable node; its content build is deferred, never boxed out.}
\fnote{\textbf{Deliverable A2 ($\rorder$, secondary).} Land the combinatorial $\rorder=a(\ell-a)+1$ fully
(same $(\ell,a)$ data as $\rlct$, a clean count). The \emph{analytic} identity ``combinatorial count $=$
pole multiplicity'' is the named at-risk seam (brief standing decision~6): it rides inside the S2 interface
for $\rlct$, but as an independent statement for $\rorder$ it needs the order-tracking form of the monomial
rule. Flag, do not hide --- $\rorder$ is the second deliverable, driven only as far as it cleanly reaches.
The reproduction of Lemmas 4--5 here transcribes the two partial-sum envelopes, the two-condition Lemma~4
characterisation, and the two-sided (upper-union $+$ explicit-construction) Lemma~5 count
(p.25--26, red-team-verified \texttt{verify-repro-s4s5.md} after an initial single-family/single-range
prose was caught and corrected); the full Case-1(2) increment bookkeeping is referenced to the images, as
$\rorder$ is off the $\rlct$ critical path.}

% =====================================================================
\section{The $L=2$ / reduced-rank-regression case, fully worked}
\label{sec:rrr}
% =====================================================================

Theorem~\ref{thm:rrr} (RRR) is Theorem~\ref{thm:main} specialised to $L=2$ --- the paper's benchmark and
the expedition's Deliverable~3 (a self-contained Lean model $+$ theorem, later exposed as the general
theorem's $L=2$ instance). The formula template is \emph{identical} once $(\ell,a,\Ms,\mathcal M)$ are
fixed: $\text{Thm 1 expression}-\text{Thm 2 expression}|_{L+1=3}=0$ exactly
(\texttt{verify-arith-groundtruth.md} Check~4). So the only content beyond Theorem~2 is the $L=2$
selection's totality and the two explicit regimes.

\subsection{The selection is total and explicit at $L=2$}

With three reduced widths $M^{(1)},M^{(2)},M^{(3)}$, Definition~\ref{def:Mset-L2} is a four-way split:
$\mathcal M=\{1,2,3\}$ (all three, $\ell=2$) when no width reaches half the sum; otherwise drop the one
\emph{large} width $M^{(s)}\ge$ (sum of the other two), leaving a size-2 set ($\ell=1$). Unlike general
Def.~3, this $L=2$ selector \emph{picks a set} wherever a width dominates --- the total version general
Def.~3 fails to be (the (T-D)/(F-1) gap). It is still not total under ties among the two smallest widths
(decided only by listing order), which is why even at $L=2$ the geometric $\Mval_{\min}$ is the clean
primitive.

\subsection{The two regimes (closed form)}

Write the Thm-3 prefactor $\rho_r:=\tfrac{-r^2+r(H^{(1)}+H^{(3)})}{2}$.

\emph{(a) Balanced ($\ell=2$, $\mathcal M=\{1,2,3\}$).} $P=M^{(1)}+M^{(2)}+M^{(3)}$, $\Ms=\lceil P/2\rceil$,
$a=P-2(\Ms-1)\in\{1,2\}$, and
\[
 \rlct=\rho_r+\frac{a(2-a)}{8}-\frac14\Big(\frac P2\Big)^2
       +\frac12\big(M^{(1)}M^{(2)}+M^{(1)}M^{(3)}+M^{(2)}M^{(3)}\big),\qquad \rorder=a(2-a)+1.
\]
Equivalently $\rlct=\rho_r+\tfrac12\Mval_{\min}$ with
$2\rlct_{\mathrm{core}}=\tfrac12(q_1^2+q_2^2-m_1^2-m_2^2-m_3^2)$, $m$ the three widths, $q$ the balanced
$2$-split of $P$.

\emph{(b) Dominated ($\ell=1$; e.g.\ $M^{(3)}\ge M^{(1)}+M^{(2)}$, $\mathcal M=\{1,2\}$).} Then $\Ms=P$,
$a=1$, $a(\ell-a)=0$ and $\ell-1=0$ kill the middle terms, leaving the bare product of the two smaller
widths:
\[
 \rlct=\rho_r+\tfrac12\,M^{(1)}M^{(2)},\qquad \rorder=1,
\]
matching the clean form $2\rlct_{\mathrm{core}}=\tfrac12\big((m_1+m_2)^2-m_1^2-m_2^2\big)=m_1m_2$.

\subsection{Ground-truth cross-checks (all PASS, \texttt{verify-arith-groundtruth.md})}

\[
\begin{array}{l|c|c|c|c|c|c}
 (H),r & \mathcal M & \ell & \Ms & a & \rlct & \rorder\\\hline
 (2,2,2),\,r{=}0 & \{1,2,3\} & 2 & 3 & 2 & 3/2 & 1\\
 (2,1,2),\,r{=}0 & \{1,2,3\} & 2 & 3 & 1 & 1 & 2\\
 (1,1,3),\,r{=}0 & \{1,2\}\ (\text{drop }M^{(3)}) & 1 & 2 & 1 & 1/2 & 1
\end{array}
\]
(For $(2,2,2,2)$, the $L=3$ equal-width case, the same closed form gives $\rlct=3/2$, $\rorder=3$ --- the
equal-width Example of \S\ref{def:Mset}, also PASS.)

\fnote{\textbf{Deliverable 3 (RRR), and a design caveat.} (a) Define the RRR model
($\mathrm{Params}$ for $L=2$, $\mathrm{dlnLoss}$, $\mathtt{rlctRRR}$) and prove Theorem~\ref{thm:rrr} via
the $L=2$ resolution the build already has. (b) Once general-$L$ lands, expose Theorem~\ref{thm:rrr} as
its $L=2$ instance. \emph{Caveat (operator, 2026-06-23):} this reproduction is a faithful reference, not a
binding Lean-API spec --- the Lean should be designed \textbf{general-$L$-first} with RRR as the exposed
instance, not RRR-first; do not let the $L=2$ closed-form split shape the general \texttt{aoyagiLambda}
API. The RRR regimes above are the cross-check ground truth, not the data model.}

% =====================================================================
\section*{Running ledger of typos / underspecifications (to collect)}
% =====================================================================
\begin{enumerate}
\item[(T-D)] \textbf{Definition 3 condition (iii) sign --- a genuine paper typo}
(image-verified p.08; \texttt{verify-def3} 2nd pass corrected the extracted-text and our
first transcription, which both had ``$\ge$''). The paper \emph{prints}
$\sum_k M^{(S_k)}\le(\ell-1)M^{(s)}$ for non-members. As printed (``$\le$'') it contradicts
the paper's own Theorem~1: at $\ell=1$ it forces $\sum_k M^{(S_k)}\le0$ --- impossible for
positive widths --- so it rejects every size-2 $\mathcal M$, exactly the sets Thm~1 uses
(witness $M=[1,1,3]$: Thm~1 picks $\{1,1\}$; printed Def.~3 rejects it), and it is
\emph{empty on the majority}: $105/216$, $476/625$, $860/1024$ at $L{=}2/3/4$. The intended
sign is almost certainly ``$\ge$'' (under which max-$\ell$ Def.~3 reduces to Thm~1's $L=2$
selector on all $512$ tested, and is never empty). \emph{This $\le$ is the kind of defect
the expedition exists to surface.}
\item[(F-1)] \textbf{Def.~3 is not a clean primitive either way.} Printed ``$\le$'':
empty / Thm-1-contradictory (T-D). Intended ``$\ge$'': total-where-nonempty but
\emph{ambiguous in $\ell$} (several prefixes admissible; ``$\ell+1$ smallest'' does not pin
$\ell$ without the max-$\ell$ reading). \textbf{Fix (verified sound,
\texttt{verify-def3}+\texttt{verify-arith}).} Take
$\rlct_{\mathrm{core}}:=\tfrac12\Mval_{\min}$, $\Mval_{\min}=\min_t\Mval(t)$ over the
admissible nested-rank lattice with \textbf{$t_L=0$} (load-bearing --- else
$\Mval_{\min}\equiv0$; structural totality: $t{=}0$ admissible, $\Mval(0)=M^{(1)}M^{(2)}<\infty$).
It is permutation-invariant (0 violations --- licenses sorting), equals the clean form
$\tfrac12(\sum_i q_i^2-\sum_k m_k^2)$, and \emph{matches the paper's total $L=2$ RLCT on all
$216$ vectors, including the $105$ where printed Def.~3 is empty} --- the geometric primitive
is right exactly where Def.~3 fails. Def.~3 (under ``$\ge$'', max-$\ell$) is a computed
consequence on its valid locus; the $L=2$ fixed Def.~\ref{def:Mset-L2} is already total.
\fnote{Lean \texttt{aoyagiLambda} $=\tfrac{-r^2+r(H^{(1)}+H^{(L+1)})}{2}+\tfrac12\Mval_{\min}$,
$\Mval_{\min}$ a $\min$ over the \textbf{$t_L=0$} admissible-profile lattice; sort widths;
printed form proven equal where Def.~3 applies.}
\qmark{\textbf{Still-owed theorem:} $\tfrac12\Mval_{\min}=\text{Thm-2 core}$ is verified by
finite exact enumeration ($L\le6$, bounded widths), \emph{not} proven $\forall(L,M)$. The
general-$L$ proof IS $\codim S(t)=\Mval(t)$ + the resolution (\S\ref{ssec:blowup}) --- the
mountain (R1, \#18). Totality, the predicate, perm-invariance are structural/derived; the
value-agreement is not.}
\item[(T-A)] \textbf{Theorem 2 rewrite (B)} (\S2.2): the middle term must be
$\big(\Ms+\tfrac{a-\ell}{\ell}\big)^2$, not $\big(\tfrac{\Ms+a-\ell}{\ell}\big)^2$;
$\Ms$ is not divided by $\ell$. (Our own transcription trap, caught + fixed in the
Thm~2 statement.)
\item[(N-1)] \textbf{Lemma 3 --- NOT a typo} (image-confirmed p.24; \texttt{verify-arith}
re-check against the image \emph{retracted} an earlier ``off-by-$\ell^2$'' typo claim
that came from the garbled extracted text). The print divides \emph{both} sides by
$\ell^2$, so $A(b):=\ell^2 b^2+\ell^2(1-2a)b+a^2\ell(\ell-1)$ is the bracket and the
printed $\min\{A\mid b=0,\dots,\ell-1\}=A(a-1)=A(a)=a\ell(\ell-a)$ is exactly right
($\partial_b A=\ell^2(1+2b-2a)$, sympy-confirmed). $\rlct$'s first variable term is
$\tfrac14\min(A/\ell^2)=\tfrac{a(\ell-a)}{4\ell}$ --- consistent.
\fnote{Formalisation note only (no typo): define $A$ as the bracket and carry
$/\ell^2$ explicitly; guard $A(a-1)$ with $a\ge1$ (invalid at $a=0$).
\emph{Process lesson:} the false ``typo'' was a text-extraction artifact --- verify
every formula against the page image, never the \texttt{.txt}.}
\item[(T-C)] \textbf{Case-2 exponent --- PINNED (elder image pass 2026-07-21).}
$M'_{S,J+1}=(M(S)-J)(M^{(S+1)}-J)$ confirmed against the p.20 image (and the Case-1
exponents against pp.16--17); non-bindingness at the flagged instance battery-verified
(\texttt{g-monument-mval-instances.py}: ``(2,2,3,2) defect non-binding'' PASS).
\item[(T-E)] \textbf{Case-2 head-reset label is RAW-width --- image-CONFIRMED
(2026-07-21), ideal-route immune.} p.20 prints $t^{(i)}_{S,J+1}=M^{(i+1)}$ (raw), Case
1(2) at p.17 prints $t^{(i)}_{S,J+1}=t^{(i)}_{s,k}$ (inherited) --- an asymmetry that at
non-monotone widths contradicts the paper's own p.22 formula
(\texttt{verify-case2-rawwidth-defect.md}). The Lean record carries no label field
(only \texttt{bexp}/\texttt{jac}); running-min governs the accumulation; the defect has
no formal representation. The formerly-open page-image residual is CLOSED.
\item[(T-F)] \textbf{T-profile totality claim is FALSE --- image-verified (2026-07-21).}
The inductive statement's display ``$T_{s,k}\le T_{s',k'}$ or $T_{s,k}\ge T_{s',k'}$''
(p.15, re-derived pp.16/17/20) fails at widths $(2,2,1,1)$: profiles $(1,1,1)$ and
$(2,1,0)$ are incomparable, both $\Mval=1$ (\texttt{g-monument-mval-instances.py}).
Nothing downstream needs totality (the $\min$ is over a set; the $b$-chain is ordered by
construction and survives --- threads 27/28); the record stores no profiles. Claim-site
notes at \S\ref{ssec:blowup} and \S\ref{sec:minimisation}.
\item[] \emph{(further items as they surface; the ranked owed-math list lives in
\texttt{verify-owed-math-audit.md})}
\end{enumerate}

\medskip\noindent\textbf{Verification status (updated 2026-07-21, elder pass).} The
\emph{formula level} (Thm~1=Thm~2$|_{L=2}$, the three Thm~2 rewrites, the equal-width
example vs.\ ground truth, the Thm~3 prefactor, Lemma~3) is \textbf{red-team PASS}
(\texttt{verify-arith}, sympy + decorrelated Codex). The \emph{page-image level} is now
closed for every load-bearing site: pp.14--22 image-verified this pass (Thm~4 statement,
inductive statement + $b_i$ + Jacobian ledger, Case-1/2 exponents and labels, the
$M_{s,k}$ formula and $\tilde t=0$ read-off), p.08 (T-D) and p.24 (N-1) earlier;
pp.25--26 red-team-verified (\texttt{verify-repro-s4s5.md}); only the $\rorder$-side
Case-1(2) increment fine-bookkeeping remains image-referenced (deferred with E). The
\emph{level lift} is now largely IN LEAN, sorry-free at the expedition tip: Lemma~1
two-sided$+$weighted (A), the guarded boxed rule (C), the QIP/codim bridge (D), the
deepest-point$+$flatten reduction for the $B=0$ corollary (GlobalHomog), the per-chart
value and negative guards (B, partial). \textbf{Remaining owed mathematics}: exactly two
sorries --- the min-over-charts change-of-variables and the coupled-atlas existence
(the monument) --- plus the consciously-deferred items ranked in
\texttt{verify-owed-math-audit.md}.

\end{document}
